\begin{trapbox}
\textbf{易错点一}
\begin{description}[style=nextline, leftmargin=2.8em, labelsep=0.5em, itemsep=0.35em]
    \item[\textbf{错误想法：}] 认为乘积形式 $[f(x)-M][f(x)-N]<0$ 对任意 $M,N$ 都等价于 $N<f(x)<M$。
    \item[\textbf{正确理解：}] 该等价性前提是 $M>N$；若 $M<N$，则乘积形式对应 $M<f(x)<N$。
    \item[\textbf{错因分析：}] 忽略了条件 $M>N$ 在推导中的关键作用，默认 $M$ 是右端点、$N$ 是左端点。
\end{description}

\medskip
\textbf{易错点二}
\begin{description}[style=nextline, leftmargin=2.8em, labelsep=0.5em, itemsep=0.35em]
    \item[\textbf{错误想法：}] 在分式形式 $\frac{f(x)-N}{M-f(x)}>0$ 中随意乘以分母 $M-f(x)$。
    \item[\textbf{正确理解：}] 分母 $M-f(x)$ 符号未知，直接乘除可能翻转不等号；应利用分子分母同号直接推得 $f(x)$ 的范围。
    \item[\textbf{错因分析：}] 未意识到分式形式本身已隐含 $M-f(x)\neq0$，且不等号方向依赖于分母正负。
\end{description}

\medskip
\textbf{易错点三}
\begin{description}[style=nextline, leftmargin=2.8em, labelsep=0.5em, itemsep=0.35em]
    \item[\textbf{错误想法：}] 认为倒数形式 $\frac{1}{f(x)-N} > \frac{1}{M-N}$ 在任何情况下都成立。
    \item[\textbf{正确理解：}] 该形式仅当 $f(x)-N > 0$ 且 $M-N>0$ 时成立，否则倒数函数单调性改变。
    \item[\textbf{错因分析：}] 忽略了倒数不等式成立的前提是两边均为正数，且函数 $\frac{1}{t}$ 在正区间单调递减。
\end{description}
\end{trapbox}
